\begin{proof}[Proof of \cref{obj:lem:reciprocal-best-approximation}]
\begin{enumerate}
\item Fix an integer \(K\ge 2\) and set
\[
H:=K^2>1.
\]
We first establish the approximation formula on a general positive interval. Let
\[
0<a_{\mathrm{app}}<b_{\mathrm{app}},
\qquad m_{\mathrm{app}}\in\mathbb N_0,
\]
and define
\[
\Delta_{\mathrm{app}}:=b_{\mathrm{app}}-a_{\mathrm{app}},
\qquad
v_{\mathrm{app}}:=\frac{a_{\mathrm{app}}+b_{\mathrm{app}}}
{\Delta_{\mathrm{app}}},
\qquad
s_{\mathrm{app}}:=\sqrt{v_{\mathrm{app}}^2-1},
\]
\[
\rho_{\mathrm{app}}:=v_{\mathrm{app}}-s_{\mathrm{app}},
\qquad
E_{\mathrm{app}}:=
\frac{2\rho_{\mathrm{app}}^{m_{\mathrm{app}}}}
{\Delta_{\mathrm{app}}(v_{\mathrm{app}}^2-1)}.
\]
Since \(\Delta_{\mathrm{app}}>0\) and
\(a_{\mathrm{app}}+b_{\mathrm{app}}>\Delta_{\mathrm{app}}\), we have
\(v_{\mathrm{app}}>1\). Consequently,
\[
0\le s_{\mathrm{app}}<v_{\mathrm{app}},
\qquad
s_{\mathrm{app}}^2=v_{\mathrm{app}}^2-1,
\qquad
(v_{\mathrm{app}}-s_{\mathrm{app}})
(v_{\mathrm{app}}+s_{\mathrm{app}})=1.
\]
Thus \(\rho_{\mathrm{app}}>0\), and
\(v_{\mathrm{app}}+s_{\mathrm{app}}>1\) gives
\[
0<\rho_{\mathrm{app}}<1,
\qquad
\rho_{\mathrm{app}}^{-1}=v_{\mathrm{app}}+s_{\mathrm{app}},
\qquad
\rho_{\mathrm{app}}+\rho_{\mathrm{app}}^{-1}=2v_{\mathrm{app}}.
\]
In particular,
\[
\rho_{\mathrm{app}}^2-2v_{\mathrm{app}}\rho_{\mathrm{app}}+1=0,
\qquad E_{\mathrm{app}}>0.
\]
% lean: Causalean.Mathlib.Analysis.Approximation.Chebyshev.Reciprocal.intervalDecay_add_inv

\item Define the Chebyshev polynomials, with integer indices, by
\[
T_0(t_{\mathrm{app}})=1,\qquad
T_1(t_{\mathrm{app}})=t_{\mathrm{app}},\qquad
T_{\nu+2}(t_{\mathrm{app}})
=2t_{\mathrm{app}}T_{\nu+1}(t_{\mathrm{app}})
-T_\nu(t_{\mathrm{app}})
\quad(\nu\in\mathbb N_0),
\]
\[
T_{-\nu}=T_\nu\quad(\nu\in\mathbb N_0).
\]
Here \(t_{\mathrm{app}}\in\mathbb R\). These polynomials satisfy
\[
\deg T_\nu=|\nu|,
\qquad
T_\nu(\cos\theta_{\mathrm{app}})
=\cos(\nu\theta_{\mathrm{app}})
\quad(\nu\in\mathbb Z,\ \theta_{\mathrm{app}}\in\mathbb R).
\]
For every \(\nu\in\mathbb N_0\), define the residual polynomial
\[
\mathcal R_{\mathrm{app},\nu}(t_{\mathrm{app}})
:=
\rho_{\mathrm{app}}^{-1}T_{\nu+1}(t_{\mathrm{app}})
-2T_\nu(t_{\mathrm{app}})
+\rho_{\mathrm{app}}T_{\nu-1}(t_{\mathrm{app}}).
\]
The degree bounds for its three terms give
\[
\deg\mathcal R_{\mathrm{app},\nu}\le \nu+1.
\]
The Chebyshev recurrence also gives
\[
\mathcal R_{\mathrm{app},\nu+2}(t_{\mathrm{app}})
=
2t_{\mathrm{app}}\mathcal R_{\mathrm{app},\nu+1}(t_{\mathrm{app}})
-\mathcal R_{\mathrm{app},\nu}(t_{\mathrm{app}}).
\]
We claim that
\[
\mathcal R_{\mathrm{app},\nu}(v_{\mathrm{app}})
=
\frac{2(v_{\mathrm{app}}^2-1)}
{\rho_{\mathrm{app}}^\nu}
\qquad(\nu\in\mathbb N_0).
\]
For the two initial indices, the identities from the preceding step yield
\[
\begin{aligned}
\mathcal R_{\mathrm{app},0}(v_{\mathrm{app}})
&=(\rho_{\mathrm{app}}^{-1}+\rho_{\mathrm{app}})
v_{\mathrm{app}}-2
=2(v_{\mathrm{app}}^2-1),\\
\mathcal R_{\mathrm{app},1}(v_{\mathrm{app}})
&=\rho_{\mathrm{app}}^{-1}(2v_{\mathrm{app}}^2-1)
-2v_{\mathrm{app}}+\rho_{\mathrm{app}}
=\frac{2(v_{\mathrm{app}}^2-1)}{\rho_{\mathrm{app}}}.
\end{aligned}
\]
If the claim holds at \(\nu\) and \(\nu+1\), the recurrence gives
\[
\begin{aligned}
\mathcal R_{\mathrm{app},\nu+2}(v_{\mathrm{app}})
&=
2(v_{\mathrm{app}}^2-1)
\left(
\frac{2v_{\mathrm{app}}}{\rho_{\mathrm{app}}^{\nu+1}}
-\frac{1}{\rho_{\mathrm{app}}^\nu}
\right)\\
&=
\frac{2(v_{\mathrm{app}}^2-1)}
{\rho_{\mathrm{app}}^{\nu+2}},
\end{aligned}
\]
because \(2v_{\mathrm{app}}\rho_{\mathrm{app}}
-\rho_{\mathrm{app}}^2=1\). This proves the claim by induction, including the degree-zero case.
% lean: Causalean.Mathlib.Analysis.Approximation.Chebyshev.Reciprocal.reciprocalResidualPoly_at_shape

\item Define the affine map
\[
t_{\mathrm{app}}(z):=
v_{\mathrm{app}}-\frac{2z}{\Delta_{\mathrm{app}}}
\qquad(z\in\mathbb R).
\]
Since \(\mathcal R_{\mathrm{app},m_{\mathrm{app}}}
(v_{\mathrm{app}})>0\), the polynomial
\[
1-
\frac{\mathcal R_{\mathrm{app},m_{\mathrm{app}}}
(t_{\mathrm{app}}(z))}
{\mathcal R_{\mathrm{app},m_{\mathrm{app}}}(v_{\mathrm{app}})}
\]
is well defined and vanishes at \(z=0\). It therefore determines a unique real polynomial \(p_{\mathrm{app}}\) through
\[
z\,p_{\mathrm{app}}(z)
=
1-
\frac{\mathcal R_{\mathrm{app},m_{\mathrm{app}}}
(t_{\mathrm{app}}(z))}
{\mathcal R_{\mathrm{app},m_{\mathrm{app}}}(v_{\mathrm{app}})}
\qquad(z\in\mathbb R).
\]
The numerator has degree at most \(m_{\mathrm{app}}+1\), so division by \(z\) gives
\[
\deg p_{\mathrm{app}}\le m_{\mathrm{app}}.
\]
For every \(z\ne0\), rearranging the defining identity yields
\[
z^{-1}-p_{\mathrm{app}}(z)
=
\frac{\mathcal R_{\mathrm{app},m_{\mathrm{app}}}
(t_{\mathrm{app}}(z))}
{z\,\mathcal R_{\mathrm{app},m_{\mathrm{app}}}
(v_{\mathrm{app}})}.
\]
% lean: Causalean.Mathlib.Analysis.Approximation.Chebyshev.Reciprocal.reciprocalApproxPoly_residual

\item For \(\theta_{\mathrm{app}}\in\mathbb R\), define
\[
\alpha_{\mathrm{app},+}(\theta_{\mathrm{app}})
:=\frac{(m_{\mathrm{app}}+1)\theta_{\mathrm{app}}}{2},
\qquad
\alpha_{\mathrm{app},-}(\theta_{\mathrm{app}})
:=\frac{(m_{\mathrm{app}}-1)\theta_{\mathrm{app}}}{2},
\]
\[
F_{\mathrm{app}}^{\sin}(\theta_{\mathrm{app}})
:=
\sin\alpha_{\mathrm{app},+}(\theta_{\mathrm{app}})
-\rho_{\mathrm{app}}
\sin\alpha_{\mathrm{app},-}(\theta_{\mathrm{app}}),
\]
\[
F_{\mathrm{app}}^{\cos}(\theta_{\mathrm{app}})
:=
\cos\alpha_{\mathrm{app},+}(\theta_{\mathrm{app}})
-\rho_{\mathrm{app}}
\cos\alpha_{\mathrm{app},-}(\theta_{\mathrm{app}}).
\]
The Chebyshev cosine identity gives
\[
\begin{aligned}
\mathcal R_{\mathrm{app},m_{\mathrm{app}}}
(\cos\theta_{\mathrm{app}})
={}&
\rho_{\mathrm{app}}^{-1}
\cos((m_{\mathrm{app}}+1)\theta_{\mathrm{app}})
-2\cos(m_{\mathrm{app}}\theta_{\mathrm{app}})\\
&+\rho_{\mathrm{app}}
\cos((m_{\mathrm{app}}-1)\theta_{\mathrm{app}}).
\end{aligned}
\]
To factor this expression, use
\[
\cos((m_{\mathrm{app}}\pm1)\theta_{\mathrm{app}})
=
1-2\sin^2\alpha_{\mathrm{app},\pm}(\theta_{\mathrm{app}})
=
2\cos^2\alpha_{\mathrm{app},\pm}(\theta_{\mathrm{app}})-1,
\]
\[
\begin{aligned}
\cos(m_{\mathrm{app}}\theta_{\mathrm{app}})
-\cos\theta_{\mathrm{app}}
&=-2\sin\alpha_{\mathrm{app},+}(\theta_{\mathrm{app}})
\sin\alpha_{\mathrm{app},-}(\theta_{\mathrm{app}}),\\
\cos(m_{\mathrm{app}}\theta_{\mathrm{app}})
+\cos\theta_{\mathrm{app}}
&=2\cos\alpha_{\mathrm{app},+}(\theta_{\mathrm{app}})
\cos\alpha_{\mathrm{app},-}(\theta_{\mathrm{app}}).
\end{aligned}
\]
Substitution, together with
\(\rho_{\mathrm{app}}+\rho_{\mathrm{app}}^{-1}
=2v_{\mathrm{app}}\), produces the two identities
\[
\begin{aligned}
\mathcal R_{\mathrm{app},m_{\mathrm{app}}}
(\cos\theta_{\mathrm{app}})
&=
2(v_{\mathrm{app}}-\cos\theta_{\mathrm{app}})
-\frac{2}{\rho_{\mathrm{app}}}
\bigl(F_{\mathrm{app}}^{\sin}(\theta_{\mathrm{app}})\bigr)^2,\\
\mathcal R_{\mathrm{app},m_{\mathrm{app}}}
(\cos\theta_{\mathrm{app}})
&=
-2(v_{\mathrm{app}}-\cos\theta_{\mathrm{app}})
+\frac{2}{\rho_{\mathrm{app}}}
\bigl(F_{\mathrm{app}}^{\cos}(\theta_{\mathrm{app}})\bigr)^2.
\end{aligned}
\]
For \(z\in[a_{\mathrm{app}},b_{\mathrm{app}}]\),
\[
t_{\mathrm{app}}(z)
=\frac{a_{\mathrm{app}}+b_{\mathrm{app}}-2z}
{\Delta_{\mathrm{app}}}
\in[-1,1].
\]
Choose \(\theta_{\mathrm{app}}\in[0,\pi]\) with
\(\cos\theta_{\mathrm{app}}=t_{\mathrm{app}}(z)\).
The nonnegative squares in the two identities imply
\[
-2(v_{\mathrm{app}}-t_{\mathrm{app}}(z))
\le
\mathcal R_{\mathrm{app},m_{\mathrm{app}}}(t_{\mathrm{app}}(z))
\le
2(v_{\mathrm{app}}-t_{\mathrm{app}}(z)).
\]
Hence
\[
\left|
\mathcal R_{\mathrm{app},m_{\mathrm{app}}}(t_{\mathrm{app}}(z))
\right|
\le
2(v_{\mathrm{app}}-t_{\mathrm{app}}(z))
=\frac{4z}{\Delta_{\mathrm{app}}}.
\]
% lean: Causalean.Mathlib.Analysis.Approximation.Chebyshev.Reciprocal.reciprocalResidualPoly_bound

\item Every \(z\in[a_{\mathrm{app}},b_{\mathrm{app}}]\) is positive. The residual identity, the preceding bound, and the value of the residual polynomial at \(v_{\mathrm{app}}\) therefore give
\[
\begin{aligned}
|z^{-1}-p_{\mathrm{app}}(z)|
&=
\frac{
|\mathcal R_{\mathrm{app},m_{\mathrm{app}}}(t_{\mathrm{app}}(z))|
}{
z\,\mathcal R_{\mathrm{app},m_{\mathrm{app}}}(v_{\mathrm{app}})
}\\
&\le
\frac{4z/\Delta_{\mathrm{app}}}
{z\,\mathcal R_{\mathrm{app},m_{\mathrm{app}}}(v_{\mathrm{app}})}
=
\frac{2\rho_{\mathrm{app}}^{m_{\mathrm{app}}}}
{\Delta_{\mathrm{app}}(v_{\mathrm{app}}^2-1)}
=E_{\mathrm{app}}.
\end{aligned}
\]
Taking the supremum yields
\[
\sup_{z\in[a_{\mathrm{app}},b_{\mathrm{app}}]}
|z^{-1}-p_{\mathrm{app}}(z)|
\le E_{\mathrm{app}}.
\]
% lean: Causalean.Mathlib.Analysis.Approximation.Chebyshev.Reciprocal.reciprocalApproxPoly_upper

\item We next construct alternating equality points. Define the grid
\[
\gamma_{\mathrm{app},\iota}
:=
\frac{\iota\pi}{m_{\mathrm{app}}+1}
\qquad
(\iota\in\{0,\ldots,m_{\mathrm{app}}+1\}).
\]
It is strictly increasing from \(0\) to \(\pi\). At its points,
\[
\alpha_{\mathrm{app},+}(\gamma_{\mathrm{app},\iota})
=\frac{\iota\pi}{2},
\qquad
\alpha_{\mathrm{app},-}(\gamma_{\mathrm{app},\iota})
=\frac{\iota\pi}{2}-\gamma_{\mathrm{app},\iota}.
\]
The sine and cosine subtraction formulas consequently give, for
\(\eta_{\mathrm{app}}\in\mathbb N_0\) whenever the displayed grid indices are in range,
\[
\begin{aligned}
F_{\mathrm{app}}^{\cos}(\gamma_{\mathrm{app},2\eta_{\mathrm{app}}})
&=(-1)^{\eta_{\mathrm{app}}}
\bigl(1-\rho_{\mathrm{app}}
\cos\gamma_{\mathrm{app},2\eta_{\mathrm{app}}}\bigr),\\
F_{\mathrm{app}}^{\sin}(\gamma_{\mathrm{app},2\eta_{\mathrm{app}}})
&=\rho_{\mathrm{app}}(-1)^{\eta_{\mathrm{app}}}
\sin\gamma_{\mathrm{app},2\eta_{\mathrm{app}}},\\
F_{\mathrm{app}}^{\cos}(\gamma_{\mathrm{app},2\eta_{\mathrm{app}}+1})
&=-\rho_{\mathrm{app}}(-1)^{\eta_{\mathrm{app}}}
\sin\gamma_{\mathrm{app},2\eta_{\mathrm{app}}+1},\\
F_{\mathrm{app}}^{\sin}(\gamma_{\mathrm{app},2\eta_{\mathrm{app}}+1})
&=(-1)^{\eta_{\mathrm{app}}}
\bigl(1-\rho_{\mathrm{app}}
\cos\gamma_{\mathrm{app},2\eta_{\mathrm{app}}+1}\bigr).
\end{aligned}
\]
The magnitudes controlling the left and right endpoints satisfy
\[
1-\rho_{\mathrm{app}}\cos\theta_{\mathrm{app}}
\ge1-\rho_{\mathrm{app}}>0
\quad(\theta_{\mathrm{app}}\in\mathbb R),
\qquad
\rho_{\mathrm{app}}\sin\theta_{\mathrm{app}}\ge0
\quad(\theta_{\mathrm{app}}\in[0,\pi]).
\]
Thus, on a cell with odd right index, the cosine factor has a nonzero left value and a right value of the opposite weak sign. On a cell with positive even right index, the sine factor has this property, since
\[
F_{\mathrm{app}}^{\sin}
(\gamma_{\mathrm{app},2(\eta_{\mathrm{app}}+1)})
=
-(-1)^{\eta_{\mathrm{app}}}\rho_{\mathrm{app}}
\sin\gamma_{\mathrm{app},2(\eta_{\mathrm{app}}+1)}.
\]
Continuity and the intermediate value theorem give choices satisfying
\[
\theta_{\mathrm{app},0}:=0,
\qquad
\gamma_{\mathrm{app},\iota-1}
<\theta_{\mathrm{app},\iota}
\le\gamma_{\mathrm{app},\iota}
\quad(1\le\iota\le m_{\mathrm{app}}+1),
\]
\[
\begin{cases}
F_{\mathrm{app}}^{\sin}(\theta_{\mathrm{app},\iota})=0,
&\iota\ \text{even},\\
F_{\mathrm{app}}^{\cos}(\theta_{\mathrm{app},\iota})=0,
&\iota\ \text{odd}.
\end{cases}
\]
The left endpoints are excluded because their factor values are nonzero. Also
\(F_{\mathrm{app}}^{\sin}(0)=0\). The cell bounds therefore imply
\[
0=\theta_{\mathrm{app},0}
<\theta_{\mathrm{app},1}<\cdots
<\theta_{\mathrm{app},m_{\mathrm{app}}+1}\le\pi.
\]
This construction applies also when \(m_{\mathrm{app}}=0\).
% lean: Causalean.Mathlib.Analysis.Approximation.Chebyshev.Reciprocal.exists_reciprocalAlternatingAngles

\item Define the nodes
\[
z_{\mathrm{app},\iota}:=
\frac{a_{\mathrm{app}}+b_{\mathrm{app}}
-\Delta_{\mathrm{app}}\cos\theta_{\mathrm{app},\iota}}{2}
\qquad(0\le\iota\le m_{\mathrm{app}}+1).
\]
Because cosine is strictly decreasing on \([0,\pi]\),
\[
a_{\mathrm{app}}\le z_{\mathrm{app},0}
<z_{\mathrm{app},1}<\cdots
<z_{\mathrm{app},m_{\mathrm{app}}+1}\le b_{\mathrm{app}},
\qquad
t_{\mathrm{app}}(z_{\mathrm{app},\iota})
=\cos\theta_{\mathrm{app},\iota}.
\]
For even \(\iota\), the sine-square identity has zero square; for odd \(\iota\), the cosine-square identity has zero square. Hence
\[
\mathcal R_{\mathrm{app},m_{\mathrm{app}}}
(t_{\mathrm{app}}(z_{\mathrm{app},\iota}))
=
(-1)^\iota
2\bigl(v_{\mathrm{app}}-t_{\mathrm{app}}(z_{\mathrm{app},\iota})\bigr)
=
(-1)^\iota\frac{4z_{\mathrm{app},\iota}}{\Delta_{\mathrm{app}}}.
\]
Substituting into the residual identity gives
\[
z_{\mathrm{app},\iota}^{-1}
-p_{\mathrm{app}}(z_{\mathrm{app},\iota})
=
(-1)^\iota
\frac{4}
{\Delta_{\mathrm{app}}\,
\mathcal R_{\mathrm{app},m_{\mathrm{app}}}(v_{\mathrm{app}})}
=
(-1)^\iota E_{\mathrm{app}}.
\]
% lean: Causalean.Mathlib.Analysis.Approximation.Chebyshev.Reciprocal.reciprocalApproxPoly_alternates

\item To certify optimality, define the signed Lagrange weights
\[
\beta_{\mathrm{app},\iota}:=
\prod_{\substack{0\le\jmath\le m_{\mathrm{app}}+1\\\jmath\ne\iota}}
(z_{\mathrm{app},\iota}-z_{\mathrm{app},\jmath})^{-1},
\qquad
\Lambda_{\mathrm{app}}:=
\sum_{\jmath=0}^{m_{\mathrm{app}}+1}
|\beta_{\mathrm{app},\jmath}|,
\qquad
w_{\mathrm{app},\iota}:=
\frac{\beta_{\mathrm{app},\iota}}{\Lambda_{\mathrm{app}}}
\quad(0\le\iota\le m_{\mathrm{app}}+1).
\]
Distinctness of the nodes makes every \(\beta_{\mathrm{app},\iota}\) nonzero, so \(\Lambda_{\mathrm{app}}>0\). There are exactly
\(m_{\mathrm{app}}+1-\iota\) negative factors in the product defining
\(\beta_{\mathrm{app},\iota}\). Therefore
\[
\sum_{\iota=0}^{m_{\mathrm{app}}+1}|w_{\mathrm{app},\iota}|=1,
\qquad
w_{\mathrm{app},\iota}
=
(-1)^{m_{\mathrm{app}}+1-\iota}
|w_{\mathrm{app},\iota}|.
\]
For any real polynomial \(q_{\mathrm{app}}\) of degree at most
\(m_{\mathrm{app}}\), Lagrange interpolation gives the polynomial identity
\[
q_{\mathrm{app}}(z)
=
\sum_{\iota=0}^{m_{\mathrm{app}}+1}
\beta_{\mathrm{app},\iota}\,
q_{\mathrm{app}}(z_{\mathrm{app},\iota})
\prod_{\substack{0\le\jmath\le m_{\mathrm{app}}+1\\\jmath\ne\iota}}
(z-z_{\mathrm{app},\jmath}).
\]
Comparing coefficients of \(z^{m_{\mathrm{app}}+1}\), and then dividing by
\(\Lambda_{\mathrm{app}}\), yields
\[
\sum_{\iota=0}^{m_{\mathrm{app}}+1}
\beta_{\mathrm{app},\iota}
q_{\mathrm{app}}(z_{\mathrm{app},\iota})=0,
\qquad
\sum_{\iota=0}^{m_{\mathrm{app}}+1}
w_{\mathrm{app},\iota}
q_{\mathrm{app}}(z_{\mathrm{app},\iota})=0.
\]
This cancellation applies to \(p_{\mathrm{app}}\) as well. Combining the weight signs with the alternating residuals gives
\[
\begin{aligned}
\sum_{\iota=0}^{m_{\mathrm{app}}+1}
w_{\mathrm{app},\iota}
\bigl(z_{\mathrm{app},\iota}^{-1}
-p_{\mathrm{app}}(z_{\mathrm{app},\iota})\bigr)
&=
(-1)^{m_{\mathrm{app}}+1}E_{\mathrm{app}}
\sum_{\iota=0}^{m_{\mathrm{app}}+1}
|w_{\mathrm{app},\iota}|\\
&=(-1)^{m_{\mathrm{app}}+1}E_{\mathrm{app}}.
\end{aligned}
\]
The cancellation identities consequently imply, for every such
\(q_{\mathrm{app}}\),
\[
\sum_{\iota=0}^{m_{\mathrm{app}}+1}
w_{\mathrm{app},\iota}
\bigl(z_{\mathrm{app},\iota}^{-1}
-q_{\mathrm{app}}(z_{\mathrm{app},\iota})\bigr)
=
(-1)^{m_{\mathrm{app}}+1}E_{\mathrm{app}}.
\]
Taking absolute values and using the unit absolute mass,
\[
\begin{aligned}
E_{\mathrm{app}}
&\le
\sum_{\iota=0}^{m_{\mathrm{app}}+1}
|w_{\mathrm{app},\iota}|
\left|z_{\mathrm{app},\iota}^{-1}
-q_{\mathrm{app}}(z_{\mathrm{app},\iota})\right|\\
&\le
\sup_{z\in[a_{\mathrm{app}},b_{\mathrm{app}}]}
|z^{-1}-q_{\mathrm{app}}(z)|.
\end{aligned}
\]
The residuals are continuous on this compact positive interval, so these suprema are finite. The lower bound for every admissible polynomial and the upper bound attained by \(p_{\mathrm{app}}\) establish
\[
\inf_{\substack{q_{\mathrm{app}}\ \mathrm{real\ polynomial}\\
\deg(q_{\mathrm{app}})\le m_{\mathrm{app}}}}
\sup_{z\in[a_{\mathrm{app}},b_{\mathrm{app}}]}
|z^{-1}-q_{\mathrm{app}}(z)|
=
\frac{2\rho_{\mathrm{app}}^{m_{\mathrm{app}}}}
{\Delta_{\mathrm{app}}(v_{\mathrm{app}}^2-1)}.
\]
% lean: Causalean.Mathlib.Analysis.Approximation.Chebyshev.Reciprocal.bestUniformApproxError_reciprocal

\item Specialize to
\[
a_{\mathrm{app}}=1,\qquad b_{\mathrm{app}}=H=K^2,
\qquad m_{\mathrm{app}}=K.
\]
These parameters meet the interval and degree conditions because \(K\ge2\).
Now
\[
v_{\mathrm{app}}=\frac{1+K^2}{K^2-1},
\qquad
(1+K^2)^2-(K^2-1)^2=4K^2,
\qquad
v_{\mathrm{app}}^2-1=\frac{4K^2}{(K^2-1)^2}.
\]
Both \(K\) and \(K^2-1\) are positive, so
\[
\sqrt{v_{\mathrm{app}}^2-1}
=\frac{2K}{K^2-1},
\qquad
\rho_{\mathrm{app}}
=
\frac{1+K^2-2K}{K^2-1}
=
\frac{K-1}{K+1}.
\]
Finally,
\[
\frac{2}{\Delta_{\mathrm{app}}(v_{\mathrm{app}}^2-1)}
=
\frac{K^2-1}{2K^2}
=
\frac{H-1}{2H}.
\]
Substitution into the established interval formula gives
\[
\inf_{\substack{p\ \mathrm{real\ polynomial}\\ \deg(p)\le K}}
\sup_{z\in[1,H]}|z^{-1}-p(z)|
=
\frac{H-1}{2H}
\left(\frac{K-1}{K+1}\right)^K,
\]
as required.
% lean: Causalean.Mathlib.Analysis.Approximation.Chebyshev.Reciprocal.bestUniformApproxError_reciprocal_one_sq
\end{enumerate}
\end{proof}
